\section{Problem Formulation}\label{sec:formulation}

\paragraph{Failure domain.} A set of components that can fail together because of a shared dependency, such as one API provider.

\paragraph{Selectors.} An agent has a primary tool and $n$ candidate corroborators, of which $s$ truly share the primary's failure domain. Metadata labels each corroborator as \emph{shared} or \emph{independent}.
\code{random\_k} chooses uniformly among the $n$ candidates, so it picks a shared tool with probability $s/n$.
\code{provenance\_k} chooses uniformly among candidates labelled independent, and falls back to a uniform choice among all $n$ if there are none.

\paragraph{Label quality.} We treat metadata as a classifier of ``shares the primary's domain'':
the \emph{recall} $r$ is the probability that a truly shared tool is labelled shared;
the \emph{false-sharing rate} $\varphi$ is the probability that a truly independent tool is labelled shared;
and $J = r - \varphi$, analogous to Youden's index in diagnostics.

\paragraph{Survey quantities.} Over pairs of servers that provide the same information type, $\rho$ is the share of pairs that truly share a provider, and $m$ is the share of pairs with a missing label (at least one side names no provider); $r$, $\varphi$ and $J$ are computed on labelled pairs. Metadata is measured at two levels: \textbf{L1} is what an agent sees at run time (server name, tool names and descriptions, input schema), and \textbf{L2} adds the README. Sharing is defined at two levels: the \emph{operator level} (the organisation operating the API host the tool calls by default) and the \emph{data-origin level} (the ultimate data source, when evidence exists).
